% Options for packages loaded elsewhere
\PassOptionsToPackage{unicode}{hyperref}
\PassOptionsToPackage{hyphens}{url}
%
\documentclass[
  11pt,
  a4paper,
]{article}
\usepackage{amsmath,amssymb}
\usepackage{iftex}
\ifPDFTeX
  \usepackage[T1]{fontenc}
  \usepackage[utf8]{inputenc}
  \usepackage{textcomp} % provide euro and other symbols
\else % if luatex or xetex
  \usepackage{unicode-math} % this also loads fontspec
  \defaultfontfeatures{Scale=MatchLowercase}
  \defaultfontfeatures[\rmfamily]{Ligatures=TeX,Scale=1}
\fi
\usepackage{lmodern}
\ifPDFTeX\else
  % xetex/luatex font selection
  \setmainfont[]{Noto Sans CJK SC}
  \setmonofont[]{DejaVu Sans Mono}
\fi
% Use upquote if available, for straight quotes in verbatim environments
\IfFileExists{upquote.sty}{\usepackage{upquote}}{}
\IfFileExists{microtype.sty}{% use microtype if available
  \usepackage[]{microtype}
  \UseMicrotypeSet[protrusion]{basicmath} % disable protrusion for tt fonts
}{}
\makeatletter
\@ifundefined{KOMAClassName}{% if non-KOMA class
  \IfFileExists{parskip.sty}{%
    \usepackage{parskip}
  }{% else
    \setlength{\parindent}{0pt}
    \setlength{\parskip}{6pt plus 2pt minus 1pt}}
}{% if KOMA class
  \KOMAoptions{parskip=half}}
\makeatother
\usepackage{xcolor}
\usepackage[margin=21mm]{geometry}
\usepackage{longtable,booktabs,array}
\usepackage{calc} % for calculating minipage widths
% Correct order of tables after \paragraph or \subparagraph
\usepackage{etoolbox}
\makeatletter
\patchcmd\longtable{\par}{\if@noskipsec\mbox{}\fi\par}{}{}
\makeatother
% Allow footnotes in longtable head/foot
\IfFileExists{footnotehyper.sty}{\usepackage{footnotehyper}}{\usepackage{footnote}}
\makesavenoteenv{longtable}
\setlength{\emergencystretch}{3em} % prevent overfull lines
\providecommand{\tightlist}{%
  \setlength{\itemsep}{0pt}\setlength{\parskip}{0pt}}
\setcounter{secnumdepth}{-\maxdimen} % remove section numbering
\XeTeXlinebreaklocale "zh"
\XeTeXlinebreakskip=0pt plus 1pt
\setlength{\emergencystretch}{3em}
\sloppy
\usepackage{fancyvrb}
\DefineVerbatimEnvironment{verbatim}{Verbatim}{fontsize=\scriptsize}

\renewcommand{\arraystretch}{1.2}
\linespread{1.08}
\raggedbottom
\raggedright
\ifLuaTeX
  \usepackage{selnolig}  % disable illegal ligatures
\fi
\IfFileExists{bookmark.sty}{\usepackage{bookmark}}{\usepackage{hyperref}}
\IfFileExists{xurl.sty}{\usepackage{xurl}}{} % add URL line breaks if available
\urlstyle{same}
\hypersetup{
  hidelinks,
  pdfcreator={LaTeX via pandoc}}

\author{}
\date{}

\begin{document}

\hypertarget{ux65e2ux6709ux4e0aux754cux8bc1ux660eux7684ux7cfbux7edfux5ba1ux67e5}{%
\section{既有上界证明的系统审查}\label{ux65e2ux6709ux4e0aux754cux8bc1ux660eux7684ux7cfbux7edfux5ba1ux67e5}}

日期：2026-10-18。审查基准：\href{https://github.com/iH-3/Sum-of-Subset/commit/702b78476185eb1554a1f679c1117f3b2eb8bf25}{702b78476185eb1554a1f679c1117f3b2eb8bf25}。

本次任务是检查已有证明，不是宣布新的上界。研究对象为有限实数集 \(A\)
满足 \(\operatorname{FS}(A)\subseteq4A\) 时的最大基数
\(N\)；子集和包括空集，四重和允许重复使用元素。

\hypertarget{ux5ba1ux67e5ux7ed3ux8bbaux4e0eux8bc1ux636eux8fb9ux754c}{%
\subsection{1.
审查结论与证据边界}\label{ux5ba1ux67e5ux7ed3ux8bbaux4e0eux8bc1ux636eux8fb9ux754c}}

\textbf{本次检查未发现新的、足以推翻已审查上界的致命数学错误。当前独立闭合的结论仍为
\(71\le N\le89\)。} 归档链声称的 \(N\le84\)
没有被本次检查否定，但不能仅凭其有限程序通过，就把它改标为``整条数学证明已经独立认证''。89
元端点的研究程序明确只给部分进展，不能据此写成 \(N\le88\)。

本次覆盖所有已声明上界的依赖关系和验证入口，重点复查定理适用条件、格化顺序、实际覆盖资源、分类穷尽性及精确反证程序。对较长归档证明的关键结构和终端论证作了直接阅读核对；这不是对每一条历史模型行逐条完成了新的人工推导，也不是证明助手中的形式化认证。

另外确认一处归档内部验证入口的优化模式防护不足；两个正式完整入口均已实测拒绝这种启动方式，本次结果不受影响。详见第 6.2 节。

\begin{longtable}[]{@{}
  >{\raggedright\arraybackslash}p{(\columnwidth - 4\tabcolsep) * \real{0.25}}
  >{\raggedright\arraybackslash}p{(\columnwidth - 4\tabcolsep) * \real{0.28}}
  >{\raggedright\arraybackslash}p{(\columnwidth - 4\tabcolsep) * \real{0.47}}@{}}
\toprule\noalign{}
\begin{minipage}[b]{\linewidth}\raggedright
对象
\end{minipage} & \begin{minipage}[b]{\linewidth}\raggedright
本次判断
\end{minipage} & \begin{minipage}[b]{\linewidth}\raggedright
主要依据
\end{minipage} \\
\midrule\noalign{}
\endhead
\bottomrule\noalign{}
\endlastfoot
基础 \(N\le122\)、\(N\le115\) & 未发现新的问题；115 必须采用已有剥离勘误
& 覆盖计数、剥离、原始文献条件与基础验证 \\
独立续研 \(115\to103\to100\to99\to\cdots\to89\) & 维持已独立闭合状态 &
一般推导审查；标准及 \texttt{-O} 精确重放；89 上界逐整数交叉核验 \\
归档 \(N\le86\) & 有限证书的重放与一般论证分开评价 & 原始 19
阶段及附加前提；后段格化、AP、完整平移和资源行核对 \\
归档 \(86\to85\to84\) & 全部列明有限检查通过；主要几何路径未见新致命问题
& 11 组完整入口，其中保留的 85 套件为 140 PASS、0 FAIL、0 SKIP \\
尝试 \(89\to88\) & 仅部分进展 &
\texttt{PASS\_PROGRESS\_ONLY}，尚存精确有理松弛见证 \\
精确最大值 \(N\) & 未确定 & 本次审查不缩小数值区间 \\
\end{longtable}

\hypertarget{ux7248ux672cux4f9dux8d56ux548cux5fc5ux987bux4fddux7559ux7684ux52d8ux8bef}{%
\subsection{2.
版本、依赖和必须保留的勘误}\label{ux7248ux672cux4f9dux8d56ux548cux5fc5ux987bux4fddux7559ux7684ux52d8ux8bef}}

审查使用独立工作分支，未修改旧证明、旧证书或冻结归档的内容。新增文件仅为审查说明、重放驱动、直接集合检查和结果回执。所有核验输出均由本次运行产生；源文件哈希、输出哈希和命令保存在结果包中。

逻辑主线为：先利用上一基数的排除结果将零加入假想反例；再整数化、剥离并建立精确覆盖预算；证明共同格及外层格性后，才能把空洞和整数间隙用于必要模型；最后以完整的根集合及穷尽二分连接到精确整数反证。一个端点的排除必须与此前较大基数的排除相接，不能把单个基数的不存在直接当成全局上界。

\hypertarget{ux5265ux79bbux5269ux4f59ux7c7bux5df2ux6709ux4feeux6b63ux5fc5ux987bux7ee7ux7eedux6cbfux7528}{%
\clearpage
\subsubsection{2.1
剥离剩余类：已有修正必须继续沿用}\label{ux5265ux79bbux5269ux4f59ux7c7bux5df2ux6709ux4feeux6b63ux5fc5ux987bux7ee7ux7eedux6cbfux7528}}

2026-10-14 的审查已指出旧表述把平移前后的同名剩余类混淆。令
\(S=\operatorname{FS}(F')\)，两端平移之差为 \(t>0\)，原支持的剩余类为
\(C\)，则应在 \(Q=C+\alpha\) 中计数。每个平移后的类增加至少一个点，因此

\[|(\alpha+S)\cup(\alpha+t+S)|\ge|S|+|Q|.\]

若限制和集所占类为 \(K\)，再加入 \(K\setminus Q\) 的代表，得到至少
\(|S|+|Q\cup K|\ge|S|+|R(F)|-1\)
个点。此修正保留引理结论，适用于后续使用同一剥离引理的地方。归档中保留旧措辞是来源保存问题；读者应同时阅读
\href{proof_review_20261014.md}{已有勘误}。本次不把它再次报成新发现。

\hypertarget{ux7aefux70b9ux6865ux4fddux7559ux4e25ux683cux7684ux51cfux4e00}{%
\subsubsection{2.2
端点桥：保留严格的减一}\label{ux7aefux70b9ux6865ux4fddux7559ux4e25ux683cux7684ux51cfux4e00}}

归档 foundations 已记录并吸收修正：若低端最大空洞为 \(h\)、总质量为
\(L\)，被保证填满的中央区间是 \([h+1,L-h-1]\)。加入新系数
\(x\)，两段整数区间重叠或相邻的充分条件为

\[x\le L-2h-1.\]

不能误用 \([h+1,L-h]\) 或放宽为 \(x\le L-2h\)。本次检查的 P4
和高空洞桥接使用修正后的门槛。原勘误位于
\path{verification/upper84/audit/verification_program/foundations/proof/endpoint_bridge_correction.txt}。

\hypertarget{ux539fux59cbux6587ux732eux4e0eux4e00ux822cux5f15ux7406}{%
\subsection{3.
原始文献与一般引理}\label{ux539fux59cbux6587ux732eux4e0eux4e00ux822cux5f15ux7406}}

\begin{longtable}[]{@{}
  >{\raggedright\arraybackslash}p{(\columnwidth - 2\tabcolsep) * \real{0.31}}
  >{\raggedright\arraybackslash}p{(\columnwidth - 2\tabcolsep) * \real{0.69}}@{}}
\toprule\noalign{}
\begin{minipage}[b]{\linewidth}\raggedright
输入
\end{minipage} & \begin{minipage}[b]{\linewidth}\raggedright
核对内容及结论
\end{minipage} \\
\midrule\noalign{}
\endhead
\bottomrule\noalign{}
\endlastfoot
Carpenter--Defant--Kravitz，v2 &
稳定性定理用于正实数集，\(r\ge4\)、\(0\le M\le r-4\)；小增量引理的前缀大小条件单独核对。不能把有限整数样本代替实数集结论。 \\
Balandraud--Girard--Griffiths--Hamidoune & 对照 PDF 原页核实 Theorem 1.5
的非严格不等号、对称非零集及非周期条件；不凭 OCR 错读严格号。Lemma 3.6
的反对称前提与后续 Kneser 分支分开。 \\
Stanchescu，1996 & 区分短直径 Theorem 4 与严格逆结论 Corollary
2；短直径陈述没有额外的 gcd 假设。严格逆门槛中的 \texttt{\textless{}}
不能替换为 \texttt{≤}。 \\
Lev，1998 &
直接核对作者稿：\(B\subseteq[1,L]\)、\(|B|=s\)、\(L\le3s/2-2\)
时得到中央子集和区间；不能先任意增加互素假设，也不能删掉密度门槛。 \\
\end{longtable}

本次核对的一般推导包括：非等差缺陷有限基例到无限归纳的连接；最后一个 AP
前缀的分情形；最大元加入时的新缺陷和被消除缺陷的恒等式；变阈值空洞递推及保留前缀的跨度、质量不等式。尖锐
\(\rho(X)\ge|X|-5\) 的使用范围仍为正文所证的
\(|X|\ge17\)，并非对任意大小成立。

高空洞递推特别区分固定阈值和放宽阈值的适用域。归档的放宽 profile
不能未经证明扩展到所有 \(\beta=9\)
情形；其实际缺陷参数、同步迭代和末尾系数界须按原论证使用。

\hypertarget{ux683cux5316ux4e0eux8986ux76d6ux9884ux7b97ux5ba1ux67e5}{%
\subsection{4.
格化与覆盖预算审查}\label{ux683cux5316ux4e0eux8986ux76d6ux9884ux7b97ux5ba1ux67e5}}

\hypertarget{ux4e0dux63d0ux524dux4f7fux7528ux683cux5316ux7ed3ux8bba}{%
\subsubsection{4.1
不提前使用格化结论}\label{ux4e0dux63d0ux524dux4f7fux7528ux683cux5316ux7ed3ux8bba}}

共同格存在以前，只使用支持基数、加法增长和已成立的实数不等式。AP
支持的单剩余类与两个剩余类分别处理；后一种需要半步格及奇系数分析。归档大小
89、88、87、86、85
的后段均存在独立的旧预算格化步骤，不能反过来提前使用格化后更强的 \(G\)
常数。

重点复查了：小侧边界、严格逆失败方向、半步 AP
分支、实际保留前缀的交叠、\(D_1\) 与 \(D\)
的逐层整性，以及原始核心的本原性。对于外层格性，新链只用单位系数消元；不能用系数
2 或更大整数的除法从格中推出原变量在格中。四空单元的 30
种部分格性例外仍保留弱间隙，没有被误删。

\hypertarget{ux6bcfux4efdux635fux5931ux53eaux6309ux5df2ux8bc1ux660eux7684ux8d44ux6e90ux5173ux7cfbux8ba1ux5165}{%
\subsubsection{4.2
每份损失只按已证明的资源关系计入}\label{ux6bcfux4efdux635fux5931ux53eaux6309ux5df2ux8bc1ux660eux7684ux8d44ux6e90ux5173ux7cfbux8ba1ux5165}}

对同一目标 \(F\) 的一组覆盖资源 \(W_i\)，容量亏损为

\[G=\sum_i c_i-\left|F\cap\bigcup_iW_i\right|.\]

不同资源组共享全局预算时，要求容量标签不重复；几何集合可以交叠。完整双和集平移只记一次容量。配对单元只有在证明实际贡献者唯一之后，才可以按该成员的原生损失计入。

若原并集的容量缺口下界为 \(H\)，成员的目标外损失下界为
\(\lambda_i\)，合法关系为 \(G\ge H+\lambda_i\)（每次选一个）以及
\(G\ge\sum_i\lambda_i\)。一般不能推出
\(G\ge H+\sum_i\lambda_i\)：成员在目标外可能重叠。已截断到目标的包络也不能再加该组成员的同一份截断损失。

归档 forest
论证按每一点的资源重数计数。带空顶点的扩展只在指定臂上至多一个空顶点、且相反臂的受限支持分离时使用；空资源没有恢复容量。多组
forest
若重用资源，使用分别成立的预算不等式及明确倍数，不能冒充一组资源互不重复的
forest。

本次直接集合检查覆盖完整平移包络、截断包络、核心分组、双和集删除包络，以及核心/普通/内层二分的几何行。120
个确定性构造集合共 66,615
项检查通过，并包含两个针对重复计损失的反例控制。它们验证的是行的编码；这些集合无需满足原始全局可行性条件，因此该测试本身不是基数排除证明。

\hypertarget{ux5f52ux6863-8584-ux7684ux5206ux7c7bux4e0eux7ec8ux7aefux8defux5f84}{%
\subsection{5. 归档 85、84
的分类与终端路径}\label{ux5f52ux6863-8584-ux7684ux5206ux7c7bux4e0eux7ec8ux7aefux8defux5f84}}

对照汇总报告、结构补充和 size-85
新报告，审查重点如下。表中的``核对''指所列关键路径，不表示已逐条人工重建所有历史模型行。

\begin{longtable}[]{@{}
  >{\raggedright\arraybackslash}p{(\columnwidth - 2\tabcolsep) * \real{0.16}}
  >{\raggedright\arraybackslash}p{(\columnwidth - 2\tabcolsep) * \real{0.84}}@{}}
\toprule\noalign{}
\begin{minipage}[b]{\linewidth}\raggedright
家族
\end{minipage} & \begin{minipage}[b]{\linewidth}\raggedright
已核对的关键义务
\end{minipage} \\
\midrule\noalign{}
\endhead
\bottomrule\noalign{}
\endlastfoot
结构入口 & 缺侧/消失侧剥离；相反对的奇偶数；原预算；AP 与非
AP；逆失败例外；格化之后才使用整数空洞 \\
Z0、P0 & 无空单元的容量与源包络矛盾 \\
Z1、P1 & 唯一空单元；中央区间定位；整数性；双 cap 排除与唯一剩余形式 \\
Z2、P2 & 实际活动性；C7 高端损失；两空单元穷尽；同一预算下的有根
forest；末端六边矛盾 \\
Z3 & A0 加权端点与轻/重裁切；A1
相反实际贡献者；物理目标方向与端点误差源 \\
P3 & 稠密种子到所有系数的桥接；边界空洞持续；A0/A1 分类；m1、m2、m3
的终端路径 \\
Z4 & A-E、A-H、B 保留表；F/O/S/E
卡片；两个配对单元均空时的无贡献者终端；全部 H 命中和 alpha 端点状态 \\
P4 & 先定位到 \(z=w=2\)，再证明外层整性；CORE8 保留；相同/相反贡献者与
T5；唯一 1 的位置不混同于配对贡献者 \\
Z5 & A0/A1/B 实际活动性分类；13 种保留 P 模板；两配对均空；6 种 pair1
空情形；C5 最终缺口、32 类实际尾部及精简 forest \\
P5、H &
从原始空洞域出发；中央区间与反射定位；高空洞整族的成对缺陷下界；不能预设
\(w\le z\) \\
\end{longtable}

较紧的 size-85 论证特别检查了两个顺序问题：Z5 alpha
缺口先提高固定损失，再推出端点正常性；C5
最终缺口用实际物理目标容量得到截断，再枚举实际尾部标签。不能借用较大基数的旧``正常端点''结论，也不能让非活跃配对成员获得容量。

仍未完成的是：把归档的每一种通用行族及分类分支与实际集合逐项建立新的独立对应。若只为把独立上界从 89 推进到 84，最短复核路径可直接接在新链之后：依次排除大小 89、88、87、86、85，复核这五个端点调用的通用引理、行族及穷尽分类；无须重新认证已被新链替代的全部早期端点。不能用相同的有限 PASS 代替这些一般证明义务。此限制不构成 84 证明错误的证据。

\hypertarget{ux7cbeux786eux9a8cux8bc1ux5668ux5ba1ux67e5ux4e0eux672cux6b21ux8fd0ux884c}{%
\subsection{6.
精确验证器审查与本次运行}\label{ux7cbeux786eux9a8cux8bc1ux5668ux5ba1ux67e5ux4e0eux672cux6b21ux8fd0ux884c}}

Farkas 反证在非负变量模型 \(Ax\le b\) 中使用非负整数行权重
\(y\)，精确检查 \(y^TA\ge0\) 且 \(y^Tb<0\)，从而得到
\(0\le y^TAx\le y^Tb<0\)。浮点求解器只能寻找候选，不能决定接受证书。

符号区间验证器直接构造二次式，不用三个采样点拟合。对离散区间检查端点以及必要的顶点附近整数；最新上界另逐个整数重建数值矩阵，与符号模型逐行比较。根集合从定义重新枚举并比较完整键集合，二分检查两个孩子和合法分支类型，不能只核对根数或叶数相同。

\begin{longtable}[]{@{}
  >{\raggedright\arraybackslash}p{(\columnwidth - 4\tabcolsep) * \real{0.22}}
  >{\raggedright\arraybackslash}p{(\columnwidth - 4\tabcolsep) * \real{0.43}}
  >{\raggedright\arraybackslash}p{(\columnwidth - 4\tabcolsep) * \real{0.35}}@{}}
\toprule\noalign{}
\begin{minipage}[b]{\linewidth}\raggedright
运行
\end{minipage} & \begin{minipage}[b]{\linewidth}\raggedright
范围
\end{minipage} & \begin{minipage}[b]{\linewidth}\raggedright
证据含义
\end{minipage} \\
\midrule\noalign{}
\endhead
\bottomrule\noalign{}
\endlastfoot
新链重放 & 17 个入口，各运行标准模式与 \texttt{-O}，共 34 次成功运行 &
输出逐字节一致；15 个有对应已发布结果的入口亦与原结果逐字节一致 \\
最新 upper89 & 30,396 个结构组、440,904 个整数根、441,243
个最终系统；1,534 个证书向量 & 全部整数矩阵交叉检查；339
个二分；最小规范化严格余量 \(1/748\) \\
直接集合几何检查 & 120 个集合，66,615 项，固定种子 2026101802 &
有限补充检查，不是全域实际集合分类 \\
upper84 完整入口 & 11/11；其中 retained upper85 为 140/140 &
列明有限义务全部通过；程序自身不宣称形式化整个定理 \\
upper86 归档 & 19 阶段及附加前提的全新组合重放 &
使用原始范围接口拼合完整覆盖；不是一次默认串行命令；最终状态见回执 \\
\end{longtable}

归档重放的 size-88 普通范围合计应为 1,128,280 个根：1,127,649
个普通证书和 631 个补充分支父根；size-87 为 1,615,300 个根，使用 8
个互不重叠范围、22 个补充组件和最后
assemble。组合驱动只有在所有范围、组件、覆盖关系及运行前后来源清单均通过后，才输出
\texttt{PASS\_FINITE\_COMPOSITION\_ONLY}。

\hypertarget{ux5982ux5b9eux4fddux7559ux8fd0ux884cux95eeux9898}{%
\subsubsection{6.1
如实保留运行问题}\label{ux5982ux5b9eux4fddux7559ux8fd0ux884cux95eeux9898}}

初次将基础 \texttt{verify.py} 强行放在 \texttt{-I}
下运行，因其原有同目录导入方式失败；按兼容的 \texttt{-E\ -S\ -B}
启动后，标准与 \texttt{-O}
均通过。这是审查驱动的调用方式问题，不是原数学证明失败。

第一次归档目录清单包含未属于原 ZIP 的环境残留文件；逐项与 ZIP
成员比对后，将这些文件保留在冻结目录外，再作完整性检查。没有通过修改原清单来接受额外文件。另一次串行重放主动中止，退出码
130，不计为完成；完整结论只使用后续全新组合重放。

最初的链回执还未解析嵌套在 \texttt{summary} 中的状态，因此部分
\texttt{reported\_status} 为
null。最终汇总重新从本次原始输出读取状态，明确保留最初失败回执与修复后运行，不把旧回执改写成成功。

\subsubsection{6.2 内部入口的优化模式防护不足}
冻结包内部 \path{proofs/exact88/verify_exact88.py} 第 28 行使用
\texttt{assert \_\_debug\_\_} 阻止优化模式。但 Python 的 \texttt{-O}
会删除该断言本身，同时删除其他证明检查中的断言。因此这条局部防护无效；不能把自行绕过正式入口、用 \texttt{-O} 运行内部脚本的日志当作认证。

本次将原始防护语句单独按 \texttt{optimize=1} 编译，确认没有断言检查指令；没有运行整份优化后的内部证明。正式入口 \path{verify_complete_upper86.py} 与 \path{verify_all.py} 使用显式条件检查，实测均以退出码 2 拒绝 \texttt{-O}。本次组合驱动通过这些正式入口或具有显式防护的冻结运行器执行，均保持断言启用，故不影响本次有限重放，也不是数学反例。

新增 \path{experiments/audit_archive_entrypoints.py} 可复现此检查（参数 \texttt{--source} 指冻结包根目录，\texttt{--output} 指 JSON 输出）。冻结包保持原字节；独立使用内部入口时，应通过受保护的正式驱动。将来维护非冻结副本时，应将这种防护改成显式 \texttt{if not \_\_debug\_\_: raise ...}，不能仍用断言检查断言是否启用。

\hypertarget{ux590dux73b0ux5165ux53e3ux548cux4ea4ux4ed8ux6587ux4ef6}{%
\subsection{7.
复现入口和交付文件}\label{ux590dux73b0ux5165ux53e3ux548cux4ea4ux4ed8ux6587ux4ef6}}

需要 Python
3.10+；证明检查仅使用标准库。所有命令在仓库根目录运行。输出目录须为新目录。以下
\texttt{audit\_new} 示例会重新计算，不复用本次结果。

\begin{verbatim}
python -I -S -B experiments/audit_upper_bound_chain.py \
  --output-dir .work/audit_new/chain --optimized --jobs 4
python -I -S -B experiments/verify_upper89_sharded.py \
  --reference .work/audit_new/chain/upper89.json \
  --output .work/audit_new/upper89_exhaustive.json \
  --work-dir .work/audit_new/parts \
  --receipt .work/audit_new/exhaustive_receipt.json --jobs 4
python -I -S -B experiments/audit_upper89_geometry.py \
  --output .work/audit_new/geometry.json
python -I -S -B verification/upper84/verify84.py --full
\end{verbatim}

先运行 \texttt{python verification/unpack\_upper86.py}，从
\path{archive/upper86/} 的 28 个分片恢复并校验原 ZIP；默认解压到
\path{.work/upper86/}。依照包内说明检查来源，然后运行：

\begin{verbatim}
python -I -S -B experiments/audit_archived_upper86.py \
  --source /absolute/path/to/subset_sums_upper86_complete_verification_20261005 \
  --output-dir .work/audit_new/archive --jobs 7
\end{verbatim}

\path{experiments/collect_proof_audit.py}
是本次历史运行目录的回执打包器，不是数学验证器。它需要本次目录布局（\texttt{replay}、\texttt{baseline\_supported}、\texttt{driver\_smoke}、\texttt{archive\_compositional}
等），严格保留启动修正的历史。新的干净重放无需制造这些失败历史，直接审阅各驱动生成的回执即可。

本次回执、标准/优化一致性、逐整数检查结果、归档阶段日志与 SHA-256
清单保存于
\path{results/2026-10-18/proof_audit/}。与已发布结果逐字节一致的大型输出只重复记录哈希；新运行的归档日志和分片证据保存在
\texttt{evidence.zip}。外部论文仅保存来源说明，不把下载全文重新分发进仓库。

\hypertarget{ux539fux59cbux6765ux6e90}{%
\subsection{8. 原始来源}\label{ux539fux59cbux6765ux6e90}}

\fontsize{9}{11}\selectfont
\begin{enumerate}
\def\labelenumi{\arabic{enumi}.}
\tightlist
\item
  Carpenter, Defant, Kravitz, \emph{Sets with few subset sums},
  \href{https://arxiv.org/html/2605.05498v2}{arXiv:2605.05498v2}，Theorem
  1.1、Lemma 2.1。
\item
  Balandraud, Girard, Griffiths, Hamidoune, \emph{Subset sums in abelian
  groups}, \href{https://arxiv.org/pdf/1112.1929}{作者稿 PDF}，Theorem
  1.5、Lemma 3.6；DOI:
  \href{https://doi.org/10.1016/j.ejc.2013.05.009}{10.1016/j.ejc.2013.05.009}。
\item
  Stanchescu, \emph{On addition of two distinct sets of integers}, Acta
  Arithmetica 75 (1996),
  191--194，\href{https://matwbn.icm.edu.pl/ksiazki/aa/aa75/aa7526.pdf}{原文
  PDF}，Theorem 4、Corollary 2；DOI:
  \href{https://doi.org/10.4064/aa-75-2-191-194}{10.4064/aa-75-2-191-194}。
\item
  Lev, \emph{On consecutive subset sums}, Discrete Mathematics 187
  (1998),
  151--160，\href{https://math.haifa.ac.il/seva/Papers/coblis.dvi}{作者
  DVI}，Theorem 1；DOI:
  \href{https://doi.org/10.1016/S0012-365X(98)80006-X}{10.1016/S0012-365X(98)80006-X}。
\end{enumerate}

\end{document}
